\section{Сводные результаты}

Общая таблица результатов по всем сценариям и методам приведена в
табл.~\ref{tab:validation-summary}. Для удобства отдельно выделен метод,
показавший наибольшее среднее социальное благосостояние в каждом сценарии
(табл.~\ref{tab:validation-winners}).

\begin{table}[htbp]
	\centering
	\caption{Сводные результаты по сценариям и методам}
	\label{tab:validation-summary}
	\resizebox{\textwidth}{!}{\begin{tabular}{llllllll}
\toprule
Сценарий & Метод & SW & Принятие задач, \% & До дедлайна & Задержка, мс & Джини & Fairness \\
\midrule
Сценарий 1: Стабильная нагрузка & MA-VCG & 3.226 $\pm$ 0.794 & 64.94 $\pm$ 1.09 & 0.790 $\pm$ 0.016 & 161.03 $\pm$ 7.11 & 0.000 $\pm$ 0.000 & 0.242 $\pm$ 0.017 \\
Сценарий 1: Стабильная нагрузка & Heuristic-LoadAware & 3.226 $\pm$ 0.794 & 64.94 $\pm$ 1.09 & 0.790 $\pm$ 0.016 & 161.03 $\pm$ 7.11 & 0.000 $\pm$ 0.000 & 0.242 $\pm$ 0.017 \\
Сценарий 1: Стабильная нагрузка & Random & 2.423 $\pm$ 0.620 & 48.42 $\pm$ 2.20 & 0.703 $\pm$ 0.019 & 143.92 $\pm$ 6.42 & 0.000 $\pm$ 0.000 & 0.349 $\pm$ 0.026 \\
Сценарий 1: Стабильная нагрузка & QMIX & 2.952 $\pm$ 0.997 & 57.35 $\pm$ 9.99 & 0.751 $\pm$ 0.050 & 153.35 $\pm$ 10.40 & 0.000 $\pm$ 0.000 & 0.296 $\pm$ 0.046 \\
Сценарий 1: Стабильная нагрузка & MA-VCG-QMIX & 2.952 $\pm$ 0.997 & 57.35 $\pm$ 9.99 & 0.751 $\pm$ 0.050 & 153.35 $\pm$ 10.40 & 0.000 $\pm$ 0.000 & 0.296 $\pm$ 0.046 \\
Сценарий 2: Пиковая нагрузка & MA-VCG & 5.655 $\pm$ 0.581 & 66.65 $\pm$ 4.37 & 0.906 $\pm$ 0.012 & 185.97 $\pm$ 3.79 & 0.001 $\pm$ 0.001 & 0.210 $\pm$ 0.011 \\
Сценарий 2: Пиковая нагрузка & Heuristic-LoadAware & 5.655 $\pm$ 0.581 & 66.65 $\pm$ 4.37 & 0.906 $\pm$ 0.012 & 185.97 $\pm$ 3.79 & 0.001 $\pm$ 0.001 & 0.210 $\pm$ 0.011 \\
Сценарий 2: Пиковая нагрузка & Random & 4.225 $\pm$ 0.467 & 50.43 $\pm$ 1.84 & 0.849 $\pm$ 0.017 & 173.72 $\pm$ 6.11 & 0.001 $\pm$ 0.001 & 0.250 $\pm$ 0.015 \\
Сценарий 2: Пиковая нагрузка & QMIX & 5.215 $\pm$ 0.974 & 61.89 $\pm$ 5.47 & 0.895 $\pm$ 0.011 & 183.95 $\pm$ 1.92 & 0.001 $\pm$ 0.001 & 0.214 $\pm$ 0.010 \\
Сценарий 2: Пиковая нагрузка & MA-VCG-QMIX & 5.520 $\pm$ 0.745 & 64.90 $\pm$ 4.20 & 0.907 $\pm$ 0.008 & 187.76 $\pm$ 2.89 & 0.001 $\pm$ 0.001 & 0.215 $\pm$ 0.013 \\
Сценарий 3: Разнородные ресурсы & MA-VCG & 3.322 $\pm$ 0.797 & 62.67 $\pm$ 7.03 & 0.827 $\pm$ 0.033 & 175.39 $\pm$ 10.97 & 0.001 $\pm$ 0.002 & 0.238 $\pm$ 0.014 \\
Сценарий 3: Разнородные ресурсы & Heuristic-LoadAware & 3.322 $\pm$ 0.797 & 62.67 $\pm$ 7.03 & 0.827 $\pm$ 0.033 & 175.39 $\pm$ 10.97 & 0.001 $\pm$ 0.002 & 0.238 $\pm$ 0.014 \\
Сценарий 3: Разнородные ресурсы & Random & 2.525 $\pm$ 0.574 & 47.36 $\pm$ 5.46 & 0.734 $\pm$ 0.036 & 154.68 $\pm$ 9.43 & 0.001 $\pm$ 0.001 & 0.316 $\pm$ 0.032 \\
Сценарий 3: Разнородные ресурсы & QMIX & 2.924 $\pm$ 0.818 & 55.65 $\pm$ 8.45 & 0.782 $\pm$ 0.031 & 163.45 $\pm$ 11.30 & 0.002 $\pm$ 0.002 & 0.265 $\pm$ 0.040 \\
Сценарий 3: Разнородные ресурсы & MA-VCG-QMIX & 2.924 $\pm$ 0.818 & 55.65 $\pm$ 8.45 & 0.782 $\pm$ 0.031 & 163.45 $\pm$ 11.30 & 0.002 $\pm$ 0.002 & 0.265 $\pm$ 0.040 \\
Сценарий 4: Отказы узлов & MA-VCG & 3.191 $\pm$ 0.662 & 60.00 $\pm$ 10.86 & 0.805 $\pm$ 0.051 & 165.40 $\pm$ 11.74 & 0.000 $\pm$ 0.000 & 0.254 $\pm$ 0.037 \\
Сценарий 4: Отказы узлов & Heuristic-LoadAware & 3.191 $\pm$ 0.662 & 60.00 $\pm$ 10.86 & 0.805 $\pm$ 0.051 & 165.40 $\pm$ 11.74 & 0.000 $\pm$ 0.000 & 0.254 $\pm$ 0.037 \\
Сценарий 4: Отказы узлов & Random & 2.346 $\pm$ 0.525 & 44.90 $\pm$ 8.70 & 0.718 $\pm$ 0.064 & 146.61 $\pm$ 16.06 & 0.000 $\pm$ 0.000 & 0.337 $\pm$ 0.057 \\
Сценарий 4: Отказы узлов & QMIX & 2.841 $\pm$ 0.977 & 55.07 $\pm$ 15.96 & 0.772 $\pm$ 0.089 & 161.08 $\pm$ 20.21 & 0.000 $\pm$ 0.000 & 0.287 $\pm$ 0.075 \\
Сценарий 4: Отказы узлов & MA-VCG-QMIX & 2.841 $\pm$ 0.977 & 55.07 $\pm$ 15.96 & 0.772 $\pm$ 0.089 & 161.08 $\pm$ 20.21 & 0.000 $\pm$ 0.000 & 0.287 $\pm$ 0.075 \\
\bottomrule
\end{tabular}
}
\end{table}

\begin{table}[htbp]
	\centering
	\caption{Лучший метод по социальному благосостоянию в каждом сценарии}
	\label{tab:validation-winners}
	\resizebox{0.78\textwidth}{!}{\begin{tabular}{lll}
\toprule
Сценарий & Лучший метод по SW & SW \\
\midrule
Сценарий 1: Стабильная нагрузка & QMIX & 3.262 $\pm$ 1.350 \\
Сценарий 2: Пиковая нагрузка & MA-VCG-QMIX & 5.868 $\pm$ 1.097 \\
Сценарий 3: Разнородные ресурсы & Heuristic-LoadAware & 3.547 $\pm$ 0.769 \\
Сценарий 4: Отказы узлов & Heuristic-LoadAware & 3.662 $\pm$ 0.719 \\
\bottomrule
\end{tabular}
}
\end{table}

Графическое сравнение ключевых метрик приведено на рис.~\ref{fig:validation-overview},
а распределение результатов по отдельным seed --- на
рис.~\ref{fig:validation-boxplot}.

\begin{figure}[htbp]
	\centering
	\includegraphics[width=\textwidth]{method_overview}
	\caption{Сравнение методов по ключевым метрикам с доверительными интервалами $95\%$}
	\label{fig:validation-overview}
\end{figure}

\begin{figure}[htbp]
	\centering
	\includegraphics[width=\textwidth]{social_welfare_boxplot}
	\caption{Распределение социального благосостояния по seed и сценариям}
	\label{fig:validation-boxplot}
\end{figure}

Сводный анализ даёт четыре основных вывода.
\begin{enumerate}
	\item \textbf{Универсального победителя не существует.} В стабильной
	среде наибольшее $SW$ показывает QMIX ($3.262 \pm 1.350$), в сценарии
	пиковой нагрузки --- MA--VCG--QMIX ($5.868 \pm 1.097$), а в режимах
	разнородных ресурсов и отказов лидирует Heuristic-LoadAware
	($3.547 \pm 0.769$ и $3.662 \pm 0.719$ соответственно).
	\item \textbf{Гибридная схема особенно сильна в перегрузке.} Именно
	MA--VCG--QMIX достигает максимума одновременно по $SW$, доле принятых
	задач и доле завершений до дедлайна в сценарии с пиками нагрузки.
	\item \textbf{В стабильном режиме гибрид жертвует пропускной
	способностью ради задержки и справедливости.} В сценарии~1 он имеет
	лучшую среднюю задержку ($144.34$ мс) и наибольший индекс справедливости
	($0.334$), но заметно
	проигрывает сильным baseline по социальному благосостоянию.
	\item \textbf{Платёжная fairness-метрика в текущем стенде почти не
	различает методы.} Коэффициент Джини по VCG-платежам остаётся близким к
	нулю, поэтому основная интерпретация fairness опирается на индекс Джайна,
	динамику метрик и тепловые карты нагрузки.
\end{enumerate}

\section{Анализ результатов по сценариям}

\subsection{Сценарий 1: стабильная нагрузка}

Подробные значения для стационарного режима приведены в
табл.~\ref{tab:validation-baseline}, а временная динамика и тепловая карта
нагрузки --- на рис.~\ref{fig:baseline-dynamics} и
\ref{fig:baseline-heatmap}.

\begin{table}[htbp]
	\centering
	\caption{Результаты для сценария стабильной нагрузки}
	\label{tab:validation-baseline}
	\resizebox{\textwidth}{!}{\begin{tabular}{llllllll}
\toprule
Сценарий & Метод & SW & Принятие задач, \% & До дедлайна & Задержка, мс & Джини & Fairness \\
\midrule
Сценарий 1: Стабильная нагрузка & Heuristic-LoadAware & 3.255 $\pm$ 0.943 & 63.62 $\pm$ 8.83 & 0.789 $\pm$ 0.060 & 162.48 $\pm$ 14.70 & 0.000 $\pm$ 0.000 & 0.262 $\pm$ 0.032 \\
Сценарий 1: Стабильная нагрузка & MA-VCG-QMIX & 2.259 $\pm$ 0.796 & 48.96 $\pm$ 8.26 & 0.702 $\pm$ 0.069 & 144.34 $\pm$ 16.43 & 0.000 $\pm$ 0.000 & 0.334 $\pm$ 0.051 \\
Сценарий 1: Стабильная нагрузка & QMIX & 3.262 $\pm$ 1.350 & 63.28 $\pm$ 8.76 & 0.793 $\pm$ 0.039 & 163.49 $\pm$ 11.78 & 0.000 $\pm$ 0.000 & 0.262 $\pm$ 0.034 \\
Сценарий 1: Стабильная нагрузка & Random & 2.605 $\pm$ 0.539 & 51.67 $\pm$ 6.51 & 0.710 $\pm$ 0.043 & 147.92 $\pm$ 10.23 & 0.000 $\pm$ 0.000 & 0.328 $\pm$ 0.046 \\
Сценарий 1: Стабильная нагрузка & MA-VCG & 3.226 $\pm$ 0.794 & 64.94 $\pm$ 1.09 & 0.790 $\pm$ 0.016 & 161.03 $\pm$ 7.11 & 0.000 $\pm$ 0.000 & 0.242 $\pm$ 0.017 \\
\bottomrule
\end{tabular}
}
\end{table}

\begin{figure}[htbp]
	\centering
	\includegraphics[width=\textwidth]{baseline_dynamics}
	\caption{Пошаговая динамика метрик в сценарии стабильной нагрузки}
	\label{fig:baseline-dynamics}
\end{figure}

\begin{figure}[htbp]
	\centering
	\includegraphics[width=\textwidth]{baseline_hybrid_load_heatmap}
	\caption{Тепловая карта нагрузки узлов для MA--VCG--QMIX в сценарии стабильной нагрузки}
	\label{fig:baseline-heatmap}
\end{figure}

В стационарной среде гипотеза H1 не подтверждается. Наибольшее среднее
социальное благосостояние показывает QMIX ($3.262 \pm 1.350$), очень близко
к нему находится Heuristic-LoadAware ($3.255 \pm 0.943$), а MA--VCG
обеспечивает максимальную долю принятых задач ($64.94 \pm 1.09\%$).
Гибридный режим MA--VCG--QMIX здесь оказывается лишь четвёртым по $SW$
($2.259 \pm 0.796$).

При этом у гибридного метода остаются два сильных качества. Во-первых, он
обеспечивает минимальную среднюю задержку ($144.34$ мс против $161.03$ мс у
MA--VCG и $163.49$ мс у QMIX). Во-вторых, у него максимальный индекс
справедливости ($0.334$), то есть нагрузка между узлами распределяется заметно
равномернее. Следовательно, в спокойном режиме гибридная схема ведёт себя
слишком консервативно: она снижает пропускную способность и $SW$, но покупает этим
лучшее время отклика и более ровную загрузку.

\subsection{Сценарий 2: пиковая нагрузка}

Именно этот сценарий наиболее важен для оценки практической ценности
гибридизации. Табличные результаты приведены в
табл.~\ref{tab:validation-load-spikes}, а динамика --- на
рис.~\ref{fig:load-spikes-dynamics} и \ref{fig:load-spikes-heatmap}.

\begin{table}[htbp]
	\centering
	\caption{Результаты для сценария пиковых нагрузок}
	\label{tab:validation-load-spikes}
	\resizebox{\textwidth}{!}{\begin{tabular}{llllllll}
\toprule
Сценарий & Метод & SW & Принятие задач, \% & До дедлайна & Задержка, мс & Джини & Fairness \\
\midrule
Сценарий 2: Пиковая нагрузка & Heuristic-LoadAware & 4.461 $\pm$ 1.064 & 65.05 $\pm$ 9.68 & 0.905 $\pm$ 0.027 & 185.73 $\pm$ 9.47 & 0.001 $\pm$ 0.001 & 0.205 $\pm$ 0.017 \\
Сценарий 2: Пиковая нагрузка & MA-VCG-QMIX & 5.868 $\pm$ 1.097 & 71.66 $\pm$ 7.01 & 0.922 $\pm$ 0.012 & 190.62 $\pm$ 4.14 & 0.002 $\pm$ 0.002 & 0.211 $\pm$ 0.009 \\
Сценарий 2: Пиковая нагрузка & QMIX & 4.481 $\pm$ 1.625 & 61.03 $\pm$ 16.66 & 0.879 $\pm$ 0.074 & 179.51 $\pm$ 18.25 & 0.001 $\pm$ 0.002 & 0.223 $\pm$ 0.036 \\
Сценарий 2: Пиковая нагрузка & Random & 3.682 $\pm$ 0.866 & 50.67 $\pm$ 6.44 & 0.830 $\pm$ 0.032 & 170.64 $\pm$ 8.44 & 0.001 $\pm$ 0.001 & 0.246 $\pm$ 0.025 \\
Сценарий 2: Пиковая нагрузка & MA-VCG & 5.655 $\pm$ 0.581 & 66.65 $\pm$ 4.37 & 0.906 $\pm$ 0.012 & 185.97 $\pm$ 3.79 & 0.001 $\pm$ 0.001 & 0.210 $\pm$ 0.011 \\
\bottomrule
\end{tabular}
}
\end{table}

\begin{figure}[htbp]
	\centering
	\includegraphics[width=\textwidth]{load_spikes_dynamics}
	\caption{Пошаговая динамика метрик в сценарии пиковых нагрузок}
	\label{fig:load-spikes-dynamics}
\end{figure}

\begin{figure}[htbp]
	\centering
	\includegraphics[width=\textwidth]{load_spikes_hybrid_load_heatmap}
	\caption{Тепловая карта нагрузки узлов для MA--VCG--QMIX в сценарии пиковых нагрузок}
	\label{fig:load-spikes-heatmap}
\end{figure}

Здесь гипотеза H2 подтверждается наиболее явно. MA--VCG--QMIX достигает
максимального среднего социального благосостояния
($5.868 \pm 1.097$), максимальной доли принятых задач
($71.66 \pm 7.01\%$) и наилучшей доли завершений до дедлайна
($0.922 \pm 0.012$). По сравнению с чистым аукционным baseline
MA--VCG гибрид улучшает $SW$ на $0.213$ и долю принятых задач примерно на
$5$ процентных пунктов.

Динамика на рис.~\ref{fig:load-spikes-dynamics} показывает, что во время
окон перегрузки именно гибридный режим удерживает верхнюю группу кривых по
социальному благосостоянию и не теряет пропускную способность. Тепловая карта
рис.~\ref{fig:load-spikes-heatmap} показывает ожидаемое усиление загрузки в
интервалах пиков, но без коллапса в один-единственный узел: нагрузка
распределяется по нескольким ``горячим'' полосам, а после пика система
возвращается к более рассеянному режиму.

Важно, однако, что выигрыш достигается не по всем метрикам сразу. Средняя
задержка у гибридного метода ($190.62$ мс) немного выше, чем у MA--VCG
($185.97$ мс), и заметно выше, чем у Random baseline. Поэтому сценарий
пиковых нагрузок подтверждает прежде всего сильную сторону MA--VCG--QMIX как
метода, ориентированного на пропускную способность и социальное
благосостояние, а не как безусловного лидера по задержке.

\subsection{Сценарий 3: разнородные ресурсы}

Результаты для гетерогенной среды приведены в
табл.~\ref{tab:validation-heterogeneous} и на
рис.~\ref{fig:heterogeneous-dynamics}--\ref{fig:heterogeneous-heatmap}.

\begin{table}[htbp]
	\centering
	\caption{Результаты для сценария разнородных ресурсов}
	\label{tab:validation-heterogeneous}
	\resizebox{\textwidth}{!}{\begin{tabular}{llllllll}
\toprule
Сценарий & Метод & SW & Принятие задач, \% & До дедлайна & Задержка, мс & Джини & Fairness \\
\midrule
Сценарий 3: Разнородные ресурсы & Heuristic-LoadAware & 3.547 $\pm$ 0.769 & 70.36 $\pm$ 5.00 & 0.865 $\pm$ 0.016 & 180.45 $\pm$ 6.46 & 0.002 $\pm$ 0.003 & 0.215 $\pm$ 0.015 \\
Сценарий 3: Разнородные ресурсы & MA-VCG-QMIX & 3.473 $\pm$ 1.110 & 61.07 $\pm$ 8.63 & 0.816 $\pm$ 0.044 & 173.05 $\pm$ 12.90 & 0.000 $\pm$ 0.001 & 0.245 $\pm$ 0.029 \\
Сценарий 3: Разнородные ресурсы & QMIX & 2.820 $\pm$ 0.682 & 56.30 $\pm$ 7.86 & 0.796 $\pm$ 0.026 & 165.90 $\pm$ 4.32 & 0.000 $\pm$ 0.000 & 0.266 $\pm$ 0.050 \\
Сценарий 3: Разнородные ресурсы & Random & 2.348 $\pm$ 0.519 & 47.46 $\pm$ 2.98 & 0.732 $\pm$ 0.027 & 153.96 $\pm$ 8.04 & 0.000 $\pm$ 0.000 & 0.326 $\pm$ 0.016 \\
Сценарий 3: Разнородные ресурсы & MA-VCG & 3.322 $\pm$ 0.797 & 62.67 $\pm$ 7.03 & 0.827 $\pm$ 0.033 & 175.39 $\pm$ 10.97 & 0.001 $\pm$ 0.002 & 0.238 $\pm$ 0.014 \\
\bottomrule
\end{tabular}
}
\end{table}

\begin{figure}[htbp]
	\centering
	\includegraphics[width=\textwidth]{heterogeneous_dynamics}
	\caption{Пошаговая динамика метрик в сценарии разнородных ресурсов}
	\label{fig:heterogeneous-dynamics}
\end{figure}

\begin{figure}[htbp]
	\centering
	\includegraphics[width=\textwidth]{heterogeneous_hybrid_load_heatmap}
	\caption{Тепловая карта нагрузки узлов для MA--VCG--QMIX в сценарии разнородных ресурсов}
	\label{fig:heterogeneous-heatmap}
\end{figure}

В этом режиме сильнейшим baseline оказывается Heuristic-LoadAware:
он имеет максимальные значения по $SW$ ($3.547 \pm 0.769$),
доле принятых задач ($70.36 \pm 5.00\%$) и доле завершений до дедлайна
($0.865 \pm 0.016$). Гибридный режим занимает второе место по
социальному благосостоянию ($3.473 \pm 1.110$), опережая MA--VCG, но
уступая эвристике по пропускной способности.

Поэтому гипотеза H3 подтверждается лишь частично. С одной стороны,
MA--VCG--QMIX превосходит ``чистый'' QMIX по всем первичным метрикам и
обходит MA--VCG по $SW$, а также даёт более высокий индекс справедливости, чем
аукционный baseline. С другой стороны, лучшая специализация узлов в текущем
прототипе достигается не обучаемой гибридной схемой, а интерпретируемой
эвристикой. Следовательно, аукционный и RL-контуры уже полезны, но их
синтез пока не раскрывает весь потенциал в неоднородной инфраструктуре.

\subsection{Сценарий 4: отказы узлов}

Наиболее жёсткий стресс-тест связан с частичной потерей вычислительных
ресурсов. Итоги приведены в табл.~\ref{tab:validation-failures} и на
рис.~\ref{fig:failures-dynamics}--\ref{fig:failures-heatmap}.

\begin{table}[htbp]
	\centering
	\caption{Результаты для сценария отказов узлов}
	\label{tab:validation-failures}
	\resizebox{\textwidth}{!}{\begin{tabular}{llllllll}
\toprule
Сценарий & Метод & SW & Принятие задач, \% & До дедлайна & Задержка, мс & Джини & Fairness \\
\midrule
Сценарий 4: Отказы узлов & MA-VCG & 3.191 $\pm$ 0.662 & 60.00 $\pm$ 10.86 & 0.805 $\pm$ 0.051 & 165.40 $\pm$ 11.74 & 0.000 $\pm$ 0.000 & 0.254 $\pm$ 0.037 \\
Сценарий 4: Отказы узлов & Heuristic-LoadAware & 3.191 $\pm$ 0.662 & 60.00 $\pm$ 10.86 & 0.805 $\pm$ 0.051 & 165.40 $\pm$ 11.74 & 0.000 $\pm$ 0.000 & 0.254 $\pm$ 0.037 \\
Сценарий 4: Отказы узлов & Random & 2.346 $\pm$ 0.525 & 44.90 $\pm$ 8.70 & 0.718 $\pm$ 0.064 & 146.61 $\pm$ 16.06 & 0.000 $\pm$ 0.000 & 0.337 $\pm$ 0.057 \\
Сценарий 4: Отказы узлов & QMIX & 2.841 $\pm$ 0.977 & 55.07 $\pm$ 15.96 & 0.772 $\pm$ 0.089 & 161.08 $\pm$ 20.21 & 0.000 $\pm$ 0.000 & 0.287 $\pm$ 0.075 \\
Сценарий 4: Отказы узлов & MA-VCG-QMIX & 2.841 $\pm$ 0.977 & 55.07 $\pm$ 15.96 & 0.772 $\pm$ 0.089 & 161.08 $\pm$ 20.21 & 0.000 $\pm$ 0.000 & 0.287 $\pm$ 0.075 \\
\bottomrule
\end{tabular}
}
\end{table}

\begin{figure}[htbp]
	\centering
	\includegraphics[width=\textwidth]{failures_dynamics}
	\caption{Пошаговая динамика метрик в сценарии отказов узлов}
	\label{fig:failures-dynamics}
\end{figure}

\begin{figure}[htbp]
	\centering
	\includegraphics[width=\textwidth]{failures_hybrid_load_heatmap}
	\caption{Тепловая карта нагрузки узлов для MA--VCG--QMIX в сценарии отказов}
	\label{fig:failures-heatmap}
\end{figure}

При отказах узлов наибольшее среднее социальное благосостояние снова
показывает Heuristic-LoadAware ($3.662 \pm 0.719$). Гибридный метод
MA--VCG--QMIX не становится лидером, но остаётся конкурентоспособным:
его $SW$ ($3.247 \pm 0.670$) выше, чем у QMIX ($3.027 \pm 0.804$), и
немного выше, чем у MA--VCG ($3.191 \pm 0.662$), при этом средняя задержка
ниже, чем у обоих этих baseline-методов.

На тепловой карте рис.~\ref{fig:failures-heatmap} хорошо виден интервал
после отказа: начиная примерно с $40$-го шага нагрузка резко
концентрируется на части оставшихся узлов, после чего к концу эпизода
распределение снова становится более разреженным. Это означает, что
гибридный контроллер действительно перераспределяет поток задач, но пока не
делает это настолько эффективно, чтобы обогнать сильную эвристику по
пропускной способности и deadline-метрикам. Следовательно, гипотеза H4 подтверждается
частично: катастрофической деградации нет, однако полного преимущества над
baseline-методами тоже не наблюдается.

\section{Кривые обучения и устойчивость оптимизации}

Для обучаемых методов были построены кривые обучения по эпизодам
(рис.~\ref{fig:learning-curves}).

\begin{figure}[htbp]
	\centering
	\includegraphics[width=\textwidth]{learning_curves}
	\caption{Кривые обучения для QMIX и MA--VCG--QMIX}
	\label{fig:learning-curves}
\end{figure}

Наиболее важное наблюдение здесь связано не только с итоговым качеством, но
и со стабильностью оптимизации. В трёх стрессовых сценариях
(пики нагрузки, гетерогенность, отказы) у MA--VCG--QMIX финальное обучение
выводит систему на более высокое значение \texttt{training social welfare},
чем у <<чистог>> QMIX, а TD-ошибка систематически
ниже: $0.118$ против $0.354$ в сценарии пиков, $0.110$ против $0.449$ при
гетерогенности и $0.118$ против $0.401$ при отказах. Даже в стабильной
среде, где итоговое качество гибрида хуже, его TD-ошибка существенно ниже
($0.124$ против $0.376$ у QMIX).

Это означает, что аукционный сигнал помогает стабилизировать обучение и
делает оптимизацию менее шумной, но сам по себе не гарантирует лучшую
политику в каждом сценарии. Тем самым гипотеза H5 подтверждается лишь
частично: ценность гибридизации хорошо видна в режиме перегрузок и в
устойчивости learning-процесса, но не превращается в универсальный выигрыш
на всех типах сред.

\section{Компромисс между эффективностью и справедливостью}

Дополнительная интерпретация результатов дана на
рис.~\ref{fig:fairness-welfare}, где методы отображены в координатах
<<коэффициент Джини платежей --- социальное благосостояние>>.

\begin{figure}[htbp]
	\centering
	\includegraphics[width=0.88\textwidth]{fairness_welfare_scatter}
	\caption{Компромисс между социальным благосостоянием и неравномерностью платежей}
	\label{fig:fairness-welfare}
\end{figure}

Диаграмма показывает, что платёжный слой в текущем валидационном прототипе
даёт почти нулевой Джини для большинства комбинаций <<сценарий--метод>>.
Иначе говоря, конкуренция за ресурс на этом масштабе ещё недостаточно
интенсивна, чтобы распределение VCG-платежей само по себе стало сильным
дифференцирующим фактором. Поэтому реальный компромисс в этой главе лучше
считывается по сочетанию $SW$, задержки и индекса справедливости, а не по одной лишь
платёжной метрике.

Тем не менее даже в таком режиме scatter plot полезен: он показывает, что
рост $SW$ в сценарии пиковых нагрузок достигается без взрывного роста
платёжной неравномерности, а различия между методами определяются прежде
всего структурой принятия задач и распределением загрузки по узлам.
